% !TEX root = ../../tree_notes.tex
% =====================================================================
%  16.1  Ancestors, LCA, and the naive walk              sec:c11-idea
%  Exports: def:c11-anc  def:c11-lcadef  prop:c11-chain  prop:c11-walk
%           ex:c11-census  fig:c11-chains  lst:c11-naive
%  (No solved problem in this section -- theory + census + naive walk.)
% =====================================================================
\section{Ancestors, the LCA, and the $\bigO{h}$ Walk}\label{sec:c11-idea}

\begin{previouslybox}[What you already know -- 4 questions]
\begin{enumerate}[nosep, leftmargin=1.8em]
  \item From \chref{c09}~(Ch.~12): the rooted-view DFS \coderef{c09-dfs3}
        stamps \(\code{depth}(v)\), \(\code{par}[v]\) and \(\code{size}(v)\)
        in one pass. \emph{(Which array do you need before any LCA question
        can even be posed?)}
  \item From \chref{c10}~(Ch.~13): the subtree fold harvests pair-values at
        the LCA (\propref{c10-pairtop}). \emph{(What was the "top node"
        $\ell$ in that construction?)}
  \item From \chref{c24}~(Ch.~5): the entry-only \(\tin/\tout\) convention
        identifies subtrees as intervals. \emph{(Can you test
        "\(u\) is an ancestor of \(v\)" with just \(\tin\)?)}
  \item From \figref{c11-tree}: the shared running tree~\(T\) rooted at~1,
        with the depth/parent overlay already printed. \emph{(Which two
        root-chains meet at node~3?)}
\end{enumerate}
\end{previouslybox}

\begin{topicintro}
Three lines of definition and one picture carry the whole chapter: root
the tree, and every node drags a chain upward to the root; two chains
climb together after they meet and never separate again.  ``Where do they
meet?'' is the LCA question.  This section proves the meeting is well
defined (common ancestors form a chain -- \propref{c11-chain}), computes a
small census of LCAs on the running tree~\(T\) \emph{by hand, from the
definition} (\exref{c11-census}) -- because every fast machine in
\S\ref{sec:c11-euler}--\S\ref{sec:c11-lifting} is just this picture
compressed -- and then writes the naive walk that is already correct and
sometimes already fast enough.  We also plant the root-dependence warning
here, where it is cheapest to understand.

\textbf{After this section you can}
\begin{itemize}[nosep]
  \item state \(\lca\) from the chain-intersection definition and compute
        any pair's LCA on paper in \(\bigO{h}\), including the
        ancestor-of-each-other shortcut case; \hfill\lv{2}
  \item prove common ancestors form a chain (so ``lowest'' is unique) and
        use the characterisation \(\lca(u,v) = u \iff u \preceq v\);
        \hfill\lv{2}--\lv{3}
  \item write the equalise-then-step-together walk and explain why its
        structure is \emph{the} template both fast machines will keep;
        \hfill\lv{3}
  \item give the worst-case cost honestly (\(\bigO{h}\) per query: fine on
        a star, fatal on a path) and name when it matters (\(Q\) large).
        \hfill\lv{2}
\end{itemize}
\end{topicintro}

% ---------------------------------------------------------------
\subsection{Chains Upward}
% ---------------------------------------------------------------

\begin{definition}{Ancestor, descendant, parent chain}{c11-anc}
Fix a root~\(r\).  The \term{ancestor chain} of~\(v\) is
\(\ancestors(v) = (v = a_0,\; a_1,\; \dots,\; a_{\depth(v)} = r)\), the
nodes of the unique route \(v \to r\) in order; \(a_j\) is the
\term{\(j\)-th ancestor} of~\(v\).  We write \(u \preceq v\) when
\(u \in \ancestors(v)\) (\(u\) is an \term{ancestor} of~\(v\), \(v\) a
\term{descendant} of~\(u\)), counting \(v \preceq v\) itself;
\(\code{par}(v) = a_1\).
\end{definition}

\noindent\textbf{Read it as:} the route, not a relation to prove.
\(\preceq\) is exactly ``appears on the way up'', and the
subtree-as-interval language of Ch.~5 translates: \(u \preceq v\) iff
\(\tin(u) \le \tin(v) \le \tout(u)\).  The self-convention
(\(v \preceq v\)) is what makes the next definition not lie about pairs
like \((9, 10)\).

\begin{definition}{Lowest common ancestor}{c11-lcadef}
For nodes \(u, v\): \(\lca(u,v)\) is the deepest node of
\(\ancestors(u) \cap \ancestors(v)\) -- the common ancestor with maximum
\(\depth\).
\end{definition}

\begin{proposition}{Common ancestors form a chain}{c11-chain}
\(\ancestors(u) \cap \ancestors(v) \ne \emptyset\) (both contain~\(r\)),
and it equals \(\ancestors(\ell)\) where \(\ell\) is its deepest element.
In particular the LCA is unique, \(w\) is a common ancestor of \(u,v\) iff
\(w \preceq \lca(u,v)\), and \(\lca(u,v) = u \iff u \preceq v\).
\end{proposition}

\begin{proofbox}
\(\ancestors(u)\) is a chain by definition (one route, no branches).  Take
any two common ancestors \(w_1, w_2\) of~\(u\) and~\(v\): both lie on the
single route \(u \to r\), and nodes of one chain are linearly ordered by
depth -- so \(w_1 \preceq w_2\) or \(w_2 \preceq w_1\).  A deepest
element~\(\ell\) exists (finite chain), and every common ancestor is an
ancestor of~\(\ell\): \(\mathit{common} = \ancestors(\ell)\).  The final
equivalence: if \(u \preceq v\), then \(\ancestors(u) \subseteq
\ancestors(v)\), so the deepest common node is~\(u\); conversely
\(\lca = u\) forces \(u\) into \(\ancestors(v)\).
\end{proofbox}

\begin{insightbox}[Three vocabularies for one node]
\(\ell\) is the \emph{deepest common ancestor}; the \emph{parting point}
where the chains from~\(u\) and~\(v\) merge; the \emph{top node / harvest
point} of Ch.~13's pair decomposition (\propref{c10-pairtop}) \emph{if}
\(u \ne v\) and neither is an ancestor of the other -- the degenerate rows
of \(\exref{c11-census}\)'s census ride on~\(\ell\) itself, one arm of
length~0.
\end{insightbox}

% ---------------------------------------------------------------
\subsection{The Parting-Point Picture on T}
% ---------------------------------------------------------------

\begin{figure}[H]
\centering
\begin{tikzpicture}[diagram, scale=0.88, transform shape,
    gv/.style={circle, draw=black!60, fill=white, minimum size=6.6mm, font=\small},
    ge/.style={draw=black!40, thick},
    eU/.style={draw=primaryblue, very thick, -{Latex[length=2mm]}},
    eV/.style={draw=accentorange, very thick, -{Latex[length=2mm]}},
    hinge/.style={gv, fill=successgreen!35, draw=green!60!black, very thick}]
  \node[gv] (n1) at (0,0)     {$1$};
  \node[gv] (n2) at (-3,-1.6) {$2$};
  \node[gv] (n3) at ( 3,-1.6) {$3$};
  \node[gv] (n4) at (-5.2,-3.2) {$4$};
  \node[gv] (n5) at (-1.6,-3.2) {$5$};
  \node[gv] (n6) at ( 1.6,-3.2) {$6$};
  \node[gv] (n7) at ( 5.2,-3.2) {$7$};
  \node[hinge] (n8) at (-1.6,-4.8) {$8$};
  \node[gv] (n9) at ( 5.2,-4.8) {$9$};
  \node[gv] (n10) at ( 5.2,-6.4) {$10$};
  \draw[ge] (n1)--(n2) (n1)--(n3) (n2)--(n4) (n2)--(n5)
            (n3)--(n6) (n3)--(n7) (n5)--(n8) (n7)--(n9) (n9)--(n10);
  % chain from 8 upward
  \draw[eU] (n8) to[bend right=12] node[annot, font=\scriptsize, left, sloped]
    {$u=8$} (n5);
  \draw[eU] (n5) to[bend right=12] (n2);
  \draw[eU] (n2) to[bend right=12] (n1);
  % chain from 10 upward
  \draw[eV] (n10) to[bend left=12] node[annot, font=\scriptsize, right, sloped]
    {$v=10$} (n9);
  \draw[eV] (n9) to[bend left=12] (n7);
  \draw[eV] (n7) to[bend left=12] (n3);
  \draw[eV] (n3) to[bend left=12] (n1);
  % hinge at 1
  \node[hinge] at (n1) {$1$};
  \node[annot, font=\scriptsize, anchor=north, text width=38mm, align=center]
    at (0,-0.55) {the two chains merge at~$1$\\(depth~$0$) and share nothing below};
  \node[annot, font=\scriptsize, anchor=west, text width=38mm, align=left]
    at (3.6,-0.3) {$\lca(8,10) = 1$: deeper on \emph{neither}\\
    chain can be common -- $8 \not\preceq 10$, $10 \not\preceq 8$};
\end{tikzpicture}
\caption{The parting-point picture on~\(T\): \(\ell = 1\) is where the
chains \(8 \to 5 \to 2 \to 1\) and \(10 \to 9 \to 7 \to 3 \to 1\) first
touch; above~\(\ell\) they coincide (both empty), below it they never meet
again (else two routes would join twice -- no cycles in a tree).  The
hinge~\(1\) is the root.}\label{fig:c11-chains}
\end{figure}

% ---------------------------------------------------------------
\subsection{Computing by Hand First (the census)}
% ---------------------------------------------------------------

\begin{example}{An LCA census straight from the definition}{c11-census}
Intersect chains, take the deepest -- no machine yet, just
\figref{c11-tree}.  On~\(T\) rooted at~1:
\[
\begin{array}{c|l|l|c}
u & v & \ancestors(u) \cap \ancestors(v) & \lca(u,v) \\ \hline
8  & 10 & \{1\} & 1\\
4  & 5  & \{2, 1\} & 2\\
6  & 7  & \{3, 1\} & 3\\
8  & 9  & \{1\} & 1\\ \hline
9  & 10 & \{9, 7, 3, 1\}\;\text{and } 9 \preceq 10 & \mathbf{9}\\
5  & 5  & \ancestors(5) = \{5,2,1\} & \mathbf{5}
\end{array}
\]
The last two rows are the shortcut cases \propref{c11-chain} licenses:
when one node is an ancestor of the other the LCA \emph{is} that node (the
chains did not part -- one contains the other), and
\(\lca(u,u) = u\).  A machine that special-cases these two lines
separately is a machine that did not read the proposition.
\end{example}

\begin{notebox}[Reading \(\exref{c11-census}\) in Ch.~13]
Chapter~13's distance-DP section (\texttt{04\_distance\_dp.tex})
references this exact example by label (\(\exref{c11-census}\)) when
proving that \(\dist(s,u)+\dist(u,v)=\dist(s,v)\) iff \(u\) is on the
route \(s\)-\(v\).  The census rows \((4,5)\) and \((6,7)\) serve as
witnesses there.
\end{notebox}

\begin{warningbox}[The LCA lives on the \emph{rooted} tree -- root changes answers]
Re-root~\(T\) at~6: the chains change, and \(\lca_6(3,6) = 6\) (since
\(6 \preceq_6 3\) now), while \(\lca_1(3,6) = 3\).  Distances, by contrast,
never moved -- the \(\bigO{1}\)-query machinery of \secref{c11-paths}
routes around this on purpose.  Fix one root for a problem, use it
everywhere, and when a statement is phrased unrooted (``the junction of
the three paths'', CSES~1136-style), any fixed root is \emph{legal}: the
answers that matter come out root-independent.
\end{warningbox}

% ---------------------------------------------------------------
\subsection{The Naive Walk: Equalise, Then Step Together}
% ---------------------------------------------------------------

\begin{proposition}{The $\bigO{h}$ walk}{c11-walk}
The following returns \(\lca(u,v)\): (i) while
\(\depth(u) > \depth(v)\): \(u \leftarrow \code{par}[u]\); (ii)
symmetrically for~\(v\); (iii) while \(u \ne v\): step \emph{both} up.
Cost is at most \(\depth(u) + \depth(v) + 1\) parent moves.
\end{proposition}

\begin{proofbox}
Phase (i)--(ii) cannot change the answer: climbing past the current LCA
is impossible while depths exceed the other's (the shallower node's chain
\emph{contains} the LCA of the pair -- \propref{c11-chain} applied to the
shortened chain), and after equalisation both nodes sit at the same depth
on their chains.  In phase (iii) equal depths mean the two walk in
lockstep and can only \emph{land} on a common node simultaneously -- the
first such landing is the deepest common ancestor: any deeper one would
have been hit one step earlier.  Termination: both reach~\(r\) at latest.
\end{proofbox}

\begin{lstlisting}[caption={The walk in five lines -- the template 16.2 and 16.3 only change how far one ``step'' may jump}, label={lst:c11-naive}]
int lca_naive(int u, int v) {
    while (dep[u] > dep[v]) u = par[u];          // phase i:  raise the deeper side
    while (dep[v] > dep[u]) v = par[v];          // phase ii: until depths match
    while (u != v) { u = par[u]; v = par[v]; }   // phase iii: lockstep until first meet
    return u;
}
\end{lstlisting}

\noindent\textbf{Believe it because:} phase (iii) is
\(\exref{c11-census}\) in motion -- the two chains \emph{merge at~\(\ell\)
and coincide above it}, so walking equal depths forward, the first
agreement \emph{must} be~\(\ell\); they can never cross without agreeing.
Every fast method below keeps this two-phase shape and merely buys a
faster ``step''.

\begin{tracebox}[Walk on T: lca(8, 10)]
\(\depth(8) = 3\), \(\depth(10) = 4\).  Phase (ii): \(v\leftarrow\code{par}[10]=9\),
now \(\depth(9)=3=\depth(8)\).  Phase (iii): step both:
\[
\begin{array}{c|c|c}
\text{step} & u & v \\ \hline
0 & 8 & 9 \\
1 & 5 & 7 \\
2 & 2 & 3 \\
3 & 1 & 1 \;\checkmark
\end{array}
\]
First agreement at depth~0: \(\lca(8,10) = 1\).  Matches the census.  Cost:
\(1 + 3 = 4\) parent moves \(= \depth(8)+\depth(10)-2\cdot\depth(1)+1\)
\(= 3+4-0+1=8\).  Wait: the \(+1\) is the final loop-iteration check, so
at most \(3+4+1=8\) pointer assignments.  For a pair on the same side
\((4,5)\): phase (i) moves \(4\to\text{nothing}\) (same depth), phase (iii)
steps \(4\to2,\;5\to2\): one move.  \(h=4\), so even the worst case is
small on~\(T\).
\end{tracebox}

\begin{center}
\footnotesize
\renewcommand{\arraystretch}{1.25}
\begin{tabular}{@{}l l l@{}}
\toprule
\textbf{Tree shape} & \textbf{Worst-case walk} & \textbf{Verdict for } \(Q\) \textbf{ queries} \\
\midrule
star / bushy (\(h = \bigO{1}\)) & \(\bigO{1}\) per query & naive \emph{wins} -- no preprocessing needed \\
balanced binary (\(h = \bigO{\log n}\)) & \(\bigO{\log n}\) & naive is already contest-fast \\
path / chain (\(h = n\)) & \(\bigO{n}\) & \(Q \ge 10^{3} \Rightarrow\) TLE: bring \secref{c11-euler} or \secref{c11-lifting} \\
\bottomrule
\end{tabular}
\end{center}

\noindent So the honest first move in an LCA problem: \emph{bound the
height}.  If the statement hands you a shallow tree (or a random one,
where Ch.~8's BST analysis says \(h = \bigO{\log n}\) in expectation),
five lines of walk beat forty lines of table.  The rest of the chapter is
for adversarial paths with large~\(Q\).

% ---------------------------------------------------------------
\subsection{Variations and Twists}
% ---------------------------------------------------------------

\begin{itemize}[nosep]
  \item \textbf{Multiple LCAs at once.}  Given \(k\) nodes,
        \(\lca(v_1,v_2,\dots,v_k)\) is the LCA of the min-depth and
        max-depth nodes in any \(\tin\)-sorted order of the set.  Proof
        sketch: the intersection of \(k\) chains is a chain (each
        intersection of two chains is a chain), so the deepest common
        element exists and equals the pairwise LCA of the two extremes.
        The Summit (\chref{c11} 99 station~12) uses this. \hfill\lv{4}
  \item \textbf{Lowest common ancestor with weights.}  If edges carry a
        length, the \emph{hinge} is the same (graph structure unchanged),
        but the \emph{distance} is \(\dist(r,u)+\dist(r,v)-2\dist(r,\ell)\).
        \secref{c11-paths} makes this the definition. \hfill\lv{3}
  \item \textbf{Ancestor test in \(\bigO{1}\).}  \(u \preceq v\) iff
        \(\tin(u) \le \tin(v) \le \tout(u)\) -- this is Ch.~5's interval
        test, restated for clarity in \propref{c11-interval}
        (\secref{c11-paths}). \hfill\lv{2}
  \item \textbf{Binary-tree specialisation.}  On a \emph{binary} tree with
        parent pointers, the walk is the classic LC~236 post-order
        solution (used in Ch.~1).  \(\lv{2}\)
  \item \textbf{Forest → super-root.}  If the input is a forest with
        \(Q\) queries spanning components, add node~0 and connect all
        component roots to it; the walk and every table below extend
        verbatim.  Queries within different components return 0. \hfill\lv{3}
\end{itemize}

% ---------------------------------------------------------------
\subsection{Common Pitfalls}
% ---------------------------------------------------------------

\begin{warningbox}[Trap 1: Forgetting the self-convention]
The LCA of a node with itself is that node (\(\lca(v,v)=v\)), not its
parent.  If your walk starts by checking \(u \ne v\) and skips the loop,
you never produce this row.  Test on \(\exref{c11-census}\)'s row~6.
\end{warningbox}

\begin{warningbox}[Trap 2: Treating \(\preceq\) as strict]
Some code writes `if (isAnc(u,v)) return u;` using a \emph{strict}
ancestor test that excludes \(u=v\).  Then \(\lca(v,v)\) falls through to
the loop and the loop never fires because \(u = v\) already: either you
return the wrong node or loop forever.  Pin the convention at the first
line: \(u \preceq v\) includes \(u = v\).
\end{warningbox}

\begin{warningbox}[Trap 3: Assuming the root is fixed]
On the tree~\(T\), \(\lca_1(8,10)=1\) but re-rooting at node~3 changes it:
the chain \(8\to5\to2\to1\to3\) and \(10\to9\to7\to3\) now meet at~3
(depth~0 in the new view).  Contest problems always say ``rooted at~1''
or give a parent array.  If they don't, \emph{read the sample} -- the
hinge answers decode the implicit root.
\end{warningbox}

% ---------------------------------------------------------------
\subsection{Practice Corner}
% ---------------------------------------------------------------

\begin{practiceset}
  \practice{Lowest Common Ancestor of a Binary Tree II}{LeetCode 1644}
    {https://leetcode.com/problems/lowest-common-ancestor-of-a-binary-tree-ii/}
    {LC Medium L2}
    {premium statement; solve from this hint: it is LC~236 (Ch.~1's mixed
     set -- the pure \emph{definition} exercise: no tables, just the
     recursion) with one twist: \(p\) and~\(q\) are \emph{not guaranteed
     to exist}.  Keep Ch.~1's post-order shape -- return the node if it is
     \(p\) or~\(q\), else the side that found something, else the merge of
     both -- but count arrivals: report an LCA only if the traversal found
     \emph{both} targets; a lone sighting is not an ancestor of a missing
     node.}
\end{practiceset}

\paragraph{Back-cites.}
LC~236 (Ch.~1 mixed set M1.1, \chref{c01}) -- the binary-tree LCA is the
base case of \(\exref{c11-census}\)'s structure.
Ch.~5 (\(\chref{c24}\)) -- the \(\tin/\tout\) interval convention
re-used in \propref{c11-interval}.
Ch.~13 (\propref{c10-pairtop}) -- the ``harvest at the hinge'' identity
that \(\exref{c11-census}\) makes concrete.

% ---------------------------------------------------------------
\begin{recallbox}[Section 16.1 in 10 seconds]
Root the tree.  Every node drags a \emph{chain} up to the root.  Two
chains meet and never part again: the meeting point is \(\lca(u,v)\).
Common ancestors form a chain (\propref{c11-chain}) so the LCA is unique.
Shortcut: \(\lca(u,v)=u\) iff \(u \preceq v\) (including \(u=v\)).  The
naive walk equalises depths and steps in lockstep; cost \(\bigO{h}\), fine
if \(h = \bigO{\log n}\) or \(Q\) small, fatal on a path with \(Q=10^5\).
\selftest{
  \item State \(\lca\) in one sentence.
  \item What does \(\preceq\) include that a strict ``below'' does not?
  \item Give the census row \((8,10)\) and explain why \(\lca=1\).
  \item The walk's three phases?  Why does lockstep find the \emph{first}
        agreement?
  \item When is the naive walk fast enough?  Give \(h\) and~\(Q\) bounds.
}
\end{recallbox}
